\cotefile{29}
\scanpage{13}
1. Données de ramification sur un schéma.

Soit $S$ un schéma (préschéma dans l'ancienne terminologie), $\underline{\mathrm{Et}}_S$ le site étale de $S$. Nous appelerons \uline{donnée de ramification sur} $S$ \deleted{la donnée} un couple $\inserted{\underline{R}=}(\underline{C},\underline{s})$, où $\underline{C}$ est une catégorie fibrée sur $\underline{\mathrm{Et}}_S$ (\deleted{\unsure{en particulier}}\inserted{i.e.} c'est une catégorie $\underline{C}$ munie d'un foncteur fibrant $p:\underline{C}\quad\underline{\mathrm{Et}}_S$) et où $\underline{s}:\underline{C}\quad\underline{\mathrm{Fl}}_S$ est un foncteur cartésien, où $\underline{\mathrm{Fl}}_S$ désigne la catégorie fibrée des \deleted{\unsure{fibreux|fibrés}} schémas relatifs sur les objets $S'$ de $\underline{\mathrm{Et}}_S$. On fera attention qu'on ne suppose pas que la catégorie fibrée $\underline{C}$ soit un champ (i.e. qu'on puisse recoller morphismes et objets de $\underline{C}$). Un objet $\inserted{M}$ de $\deleted{\underline{C}(S')}$ la catégorie fibre $\underline{C}_{S'}$ ($S'\quad\mathrm{Ob}\,\underline{\mathrm{Et}}_S$) sera appelé un modèle de ramification pour $\underline{R}$, et le schéma $\underline{s}(M)$ sur $S'$ sera appelé la réalisation géométrique de ce modèle. On fera attention qu'on ne suppose pas le foncteur $\underline{s}$ fidèle ni pleinement fidèle, donc les morphismes entre modèles de ramification ne peuvent être identifiés aux morphismes entre leurs réalisations géométriques. Dansles exemples envisagés par la suite, les $\underline{s}(M)$ sont finis sur $S'$.

Considérons le foncteur composé $\underline{C}\overset{\underline{s}}{\qquad}\underline{\mathrm{Fl}}_S\overset{\mathrm{source}}{\qquad}(\mathrm{Sch})$, d'où un foncteur (le foncteur ``réalisation géométrique'')
\[
\underline{C}\qquad(\mathrm{Sch})\,.
\]
Considérons d'autre part la catégorie fibrée sur$(\mathrm{Sch})$ des schémas relatifs étales et finis (dont la catégorie fibre en $T$ est donc la catégorie des revêtements étales finis de $T$), et soit $\underline{D}$ la catégorie fibrée sur $\underline{C}$ image inverse de la catégorie fibrée précédente par le foncteur réal\unsure{'|i}sation géométrique. La fibre de cette catégorie fibrée en l'objet $M$ de $\underline{C}$ est donc la catégorie des revêtements étales finis de la réalisation géométrique $\underline{s}(M)$ de $M$. O\textsuperscript{\unsure{N|n}} appelera \uline{revêtement de} \deleted{\unsure{type}}
% archive sheet number: 12
\scanpage{14}
Appelons \uline{crible couvrant} de $\underline{C}$ toute sous-catégorie \inserted{pleine} $\underline{C}_{\unsure{o|0}}$ de $\underline{C}$ qui est un crible (i.e. pour toute flèche, l'origine appartient à $\underline{C}_o$ s'il en est de même de l'extrémité) et tel que les $\underline{p}(M)$ pour $M\quad\mathrm{Ob}\,\underline{C}_o$ recouvrent l'objet final $S$ de $\underline{\mathrm{Et}}_S$. Pour tout $\underline{C}_o$, soit $\mathrm{Rev}^{\underline{R}}_{\underline{C}_o}(S)$ la catégorie des sections cartésiennes de $\underline{D}$ au d essus de $\underline{C}_o$. Pour $\underline{C}_o$ variable\deleted{\unsure{s}}, ces catégories forment \deleted{\unsure{un} les fibres d'une}\inserted{un système inductif de} catégories\deleted{fibrées sur la catégorie ordonnée de ces $\underline{C}_o$} (NB si $\underline{C}'_o\quad\underline{C}_o$, on obtient un foncteur restriction évident $\mathrm{Rev}^{\underline{R}}_{\underline{C}_o}(S)\quad\mathrm{Rev}^{\underline{R}}_{\underline{C}'_o}(S)$), \deleted{et} dont la limite inductive est notée $\mathrm{Rev}^{\underline{R}}(S)$ et appelée \uline{catégorie des revêtements de} $S$ \uline{de type} $\underline{R}$. Signalons que dans tous les cas que nous aurons à utiliser, les foncteurs de transition $\mathrm{Rev}^{\underline{R}}_{\underline{C}_o}(S)\quad\mathrm{Rev}^{\underline{R}}_{\underline{C}'_o}(S)$ sont pleinement fidèles, donc aussi les foncteurs $\mathrm{Rev}^{\underline{R}}_{\underline{C}_o}(S)\quad\mathrm{Rev}^{\underline{R}}(S)$, qui permettent donc d'identifier (à équivalence près) les $\mathrm{Rev}^{\underline{R}}_{\underline{C}_o}(S)$ à des sous-catégories pleines de la catégorie $\mathrm{Rev}^{\underline{R}}(S)$ des revêtements detype $S$. Nous supposerons \deleted{de plus} que pour tout $\underline{C}_o$, \deleted{le foncteur composé} et tout $X\quad\mathrm{Ob}\,\mathrm{Rev}^{\underline{R}}_{\underline{C}_o}(S)$, le foncteur composé $\underline{C}_o\overset{X}{\qquad}\underline{D}\overset{\mathrm{source}}{\qquad}(\mathrm{Sch})_{/S}$, (où le deuxième foncteur est le foncteur déduit du foncteur source sur la catégorie des s chémas finis étales relatifs,) admet une limite inductive dans $(\mathrm{Sch})_{/S}$, que nous noterons $\unsure{\mathbb{r}|\underline{r}}(X)$, et que pour $C'_o\quad\underline{C}_o$, désignant par $X'$ la restriction de $X$ à $\underline{C}'_o$, l'homomorphisme canonique $r(X'_o)\quad r(X)$ est un isomorphisme. (Ces conditions seront satisfaites dans les exemples \inserted{(de généralité croissante)} 1 à 3 ci-dessous). Les différents foncteurs $r$ se r \unsure{accordent} alors et
% archive sheet number: 13
\scanpage{15}
\deleted{\uline{type} $\underline{R}$ \uline{de} $S$ une section cartésienne de $\underline{D}$ sur $\underline{C}$. Ces revêtements de type $\underline{R}$ forment une catégorie \inserted{que nous noterons $\mathrm{Rev}_{\underline{R}}(S)$.} \handwritten{(} notée suivant les notations standard $\underset{\handwritten{\longleftarrow}}{\mathrm{Lim}}\,\underline{D}/\underline{C}$\handwritten{)}. On définira la \uline{réalisation géométrique} d'un revêtement $\handwritten{\underline{\phi}}$ de type $\underline{R}$ comme étant la limite inductive (si elle existe) dans $(\mathrm{Sch}$ du foncteur composé $\underline{C}\overset{\handwritten{\underline{\phi}}}{\qquad}\underline{D}\qquad(\mathrm{Sch})_{/S}$, où le deuxième de ces foncteurs est le foncteur source \deleted{\handwritten{(}associant à tout \unsure{revêtement}}(dé\inserted{d}uit du foncteur source sur la catégorie des schémas finis étales relatifs). Si pour tout $\handwritten{\underline{\phi}}$, cette limite inductive existe,} on obtient un foncteur réalisation géométrique
\[
\deleted{\handwritten{\unsure{\phi|r}}}\quad r:\mathrm{Rev}_{\underline{R}}(S)\qquad(\mathrm{Sch})_{/S}\,.
\]
Dans les cas qui nous intéresseront par la suite, ce foncteur $r$ sera fidèle, mais pas nécessairement pleinement fidèle; il importera donc de bien distinguer entre une revêtement de type $\underline{R}$ de $S$, et sa réalisation géométrique, le premier n'étant pas déterminée par la seule connaissance du second.

Avant d'aller plus loin, donnons quelques exemples.

\uline{Exemple} 1. Donnons-nous un schéma relatif $Z$ sur $S$, muni d'un groupe fini $G$ opérant \inserted{(à droite disons)} par $S$-automorphismes. Soit $C_{\unsure{o|0}}$ la catégorie réduite à un objet $\inserted{e}$ et dont le monoïde des endomorphismes est $G$, et prenons pour $\underline{C}$ la catégorie produit $\underline{\mathrm{Et}}_S\mathrm{x}C_o$, \deleted{pour foncteur $\underline{s}$ le foncteur $(S',e)$ qui aux les objets} La donnée d'un foncteur cartésien de \deleted{Et} $\underline{C}$ dans une catégorie fibrée $\underline{C}'$ sur $\underline{\mathrm{Et}}_S$ équivaut manifestement à la donnée d'un objet de la catégorie fibre $\underline{C}'_S$ avec une opération de $G$ dessus, de sorte que $(Z,G)$ nous \unsure{fournit} bien un foncteur cartésien $\underline{s}$. On trouve aussitôt que la catégorie $\mathrm{Rev}_{\underline{R}}(S)$ est équivalente canoniquement à la \deleted{\unsure{cat}} catégorie des revêtements finis ét\unsure{q|a}les $Z'$ de $Z$, munis d'une opération (droite de $G$ \deleted{compatible} rendant le morphisme structural $Z'\quad Z$ un $G$-morphisme.
% archive sheet number: 14
\scanpage{16}
La réalisation géométrique d'un tel $Z'$ existe si et seulement si le quotient $Z'/G$ existe, \inserted{et elle s'identifie alors à $X=Z'/G$.} Il en est ainsi si toute orbite de $G$ dans $Z$ est contenue dans un ouvert affine, ce qui implique en effet qu'il en est de même de toute orbite de $G$ dans $Z'$. \deleted{Intuitivement, \inserted{et en première approximation, les réalisations} la considération les revêtements de type $\underline{R}$ correspondent à des revêtements} Dans les cas que nous avons en vue, $Z$ est un revêtement de $S$, donc la condition sur les orbites sera automatiquement satisfaite (il suffit que $Z$ soit affine sur $S$)\deleted{plus généralement que ses fibres} de plus les réalisations géométriques $X=Z'/G$ seront automatiquement finies sur $S$, \deleted{donc}\inserted{i.e.} des revêtements. En première approximation, ce seront les revêtements de $S$ qui (en \deleted{\unsure{un|ujn}} un sens convenable ...) deviennent étales au dessus du revêtement (ramifié en général) donné $Z$, ce qui donne un sens intuitif à la notion de ``donnée de ramification''. Soit maintenant $U$ un ouvert de $S$ tel que $Z\;U$ soit un revêtement principal de groupe $G$, et supposons que $Z\;U=V$ soit dense dans $Z$ (plus généralement, que toute partie ouverte et fermée de $Z$ contenant $Z\;U$ \inserted{et stable par $G$} soit égale à $Z$). Alors le foncteur $Z'\qquad Z'\;V$ de la catégorie des $G$-revêtements étales de $Z$ dans la catégorie des $G$-revétments étales de $V$ est fidèle, d'autre part le foncteur $Z'_V\qquad Z'_{\deleted{V}}/G$ de la catégorie des $G$-revêtement s étales de $V$ dans la catégorie des revêtements étales de $U$ est une équivalence de catégories, à fortiori fidèle, d'où on conclut que le foncteur \inserted{réalisation géométrique} $Z'\qquad Z'/G=X$ est\unsure{x|f} \uline{fidèle}. Suivant le yoga général\deleted{\unsure{e|m}}, \deleted{une donnée de}\inserted{la catégorie des} revêtementsde type $\underline{R}$ \deleted{du} est alors équivalente à la catégorie des ``revêtements de $S$ munis d'une structure de type $\underline{R}$'', où par ``structure de type $\underline{R}$'' sur un revêtement $X$ de $S$ on entend la donnée d'une classe (mod isomorphisme) d'un $Z'\quad\mathrm{Rev}_{\underline{R}}(S)$ et d'un isomorphisme de $\underline{r}(Z')$ ($=Z'/G$) avec $X$. (Cette
% archive sheet number: 15
\scanpage{17}
dernière remarque aurait dû figurer avant l'exemple, elle s'applique chaque fois que le foncteur $\underline{r}$ est fidèle). \deleted{Nous pouvons également} Supposons maintenant $V$ dense dans $Z$ et de plus $Z$ \uline{normal}, alors on voit aussitôt que le foncteur $Z'\qquad Z'\;V$ est pleinement fidèle, d'où \unsure{résulte} facilement qu'il en est de même du foncteur réalisation géométrique $Z'\qquad Z'/G$; donc dans ce cas, la catégorie $\mathrm{Rev}_{\underline{R}}(S)$ \deleted{\unsure{est}} estéquivalente à une sous-catégorie pleine de la catégorie des revêtements \deleted{\unsure{de|du}} normaux de $S$,\inserted{étales au dessus de $U$, et} dont toute composante \deleted{domine} irréductible domine une composante irréductible de $S$ rencontrant $U$.

\uline{Exemple} 2 C'est une simple variante de l'exemple précédent: au lieu de se donner un $S$-schéma $Z$ à groupe fini d'opérateurs $G$, on se donne un système projectif $(Z_i,G_i)_{i\;I}$ de tels objets ($I$ ensemble ordonné filtrant). \deleted{La catégorie Rev} Le système projectif de\unsure{x|s} groupes $(G_i)$ définit une catégorie fibrée $C_{\unsure{o|0}}$ sur la catégorie associée à $I$ (les fibres de $C_o$ étant les catégories à un seul objet définies par les $G_i^{\unsure{o|0|\circ}}$), et on prendra $\underline{C}=\underline{\mathrm{Et}}_S\mathrm{x}C_o$, et on définit $\underline{s}$ de façon évidente. La catégorie $\mathrm{Rev}_{\underline{R}}(S)$ est maintenant équivalente à la catégorie $\mathrm{Lim}$ des catégories $\mathrm{Rev}_{\underline{R}_i}(S)$ définies séparément par les $(Z_i,G_i)$. Un cas intéressant est celui où les foncteurs transition (qui s'interprètent de façon évidente comme des foncteurs imagesinverses de revêtements étales à opérateurs) sont pleinement fidèles, de sorte que les $\mathrm{Rev}_{\underline{R}_i}(S)$ s'identifient (à équivalence près) à des sous-catégories pleines de $\mathrm{Rev}_{\underline{R}}(S)$ qui \unsure{epuisent|épuisent} évidemment cette dernière. Il en sera ainsi si les morphismes de transition $(Z_i,G_i)\qquad(Z_j,G_j)$ sont tels que $G_i\quad G_j$ soit \unsure{surjectif} et $Z_i\quad Z_j$ induit un isomorphisme $Z_j\quad Z_i/N_{ij}\mathord{,}_{R}$ où $N_{ij}$ est le noyau de $G_i\quad G_j$ (en effet, il en résulte que pour tout revêtement \inserted{étale $Z'_j$}
% archive sheet number: 16
\scanpage{19}
de $Z_j$, on a $Z'_j\quad Z'_{i\deleted{j}}/N_{ij}$, où $Z'_i$ est l'image inverse de $Z'_j$ sur $Z_i$). Dans ce cas, le foncteur ``réalisation géométrique'' est bien défini, et induit surchaque $\mathrm{Rev}^{\underline{R}_i}(S)$ le foncteur réalisation géométrique déjà explicité dans l'exemple \unsure{1|2}. \deleted{Bien sûr, les conditions} Notons que les \unsure{criètres|crièères|critères} signalés dans l'exemple 1 pour que le foncteur réalisation géométrique soit fidèle resp. pleinement fidèle, s'étendent de façon immédiate au cas envisagé ici.

\uline{Exemple} 3. Supposons que $\underline{R}=(\underline{C},\underline{s})$ satisfasse aux conditions suivantes. Le foncteur $\underline{s}$ se factorise en $\underline{C}\overset{\underline{s}'}{\qquad}\underline{\unsure{\mathrm{Flop}|\mathrm{Faop}}}_S\qquad\underline{\mathrm{Fl}}_S$, où $\underline{\mathrm{Flop}}_S$ désigne la catégorie fibrée sur $\underline{\mathrm{Et}}_S$ des schémas relatifs à groupe fini (opérant à droite) d'opérateurs, et $\underline{\mathrm{Flop}}_S\quad\underline{\mathrm{Fl}}_S$ est le foncteur ``oubli\deleted{\unsure{r|e}} des opérations''. \deleted{Pour toute flèche $M'\quad M$ de $\underline{C}_{S'}$, la flèche \unsure{correspondante}}\inserted{Il faut préciser que (en vue de l'exemple kummérien plus} bas) nous prenons comme morphismes d'un schéma à opérateurs $(Z,G)$ sur $S$ dans un autre $(Z',G')$ les couples $(f,u)$, où $f:X\quad X'$ est un $S$-morphisme et $u:G_S\quad G'_S$ un morphisme(localement constant,\deleted{(}pas nécessairement constant !) de schémas en groupes sur $S$, de telle façon que $f$ soit compatible avec $u$ dans le s ens évident. On suppose que pour toute flèche $M'\quad M$ dans un $\underline{C}_S$, \deleted{\unsure{déf}} la flèche correspondante $\underline{s}(M')=(Z',G')\quad\underline{s}(M)=(Z,G)$ induit un isomorphisme $Z\quad Z'/H$, où $H$ est le noyau de $G'_S\quad G_S$ (qui est un sous-groupe localement constant de $G'_S$). On suppose \deleted{que tout} aussi que pour tout objet $M$ de $\mathrm{Ob}\,\underline{C}_{S'}$, posant $\underline{s}(M)=(Z,G)$, \deleted{et pour tout $g\quad G$, il existe un endomorphisme $\underline{g}$ de $M$, uniquement déterminé, tel que}\inserted{on se donne un \deleted{automorphisme} homomorphisme $G\quad\mathrm{Aut}(M)$, tel que pour tout $g\quad G$, si $\underline{g}$ est son image,} $\underline{s}'(\underline{g})$ soit le couple $(g_Z,\mathrm{int}(g))$, et que tout \inserted{tel} endomorphisme de $M$ soit localement de la forme $\underline{g}$; on \deleted{veut} suppose de plus que \deleted{la} l'homomorphisme $G\quad\mathrm{Aut}(M)$ soit \unsure{compatibde|compatible} avec images inverses, et que pour toute flèche $M'\quad M$ dans $\underline{C}_{S'}$, considérons $M'$ (resp. $M$) comme muni ud groupe $G'$ ($G$) d'opérateurs, $M'\quad M$ soit compatible avec $G'\quad G$.
% archive sheet number: 18
\scanpage{20}
Dans ce cas, le foncteur ``réalisation géométrique'' est bien défini et induit sur chaque $\mathrm{Rev}_{\underline{R}_i}(S)$ le foncteur réalisation géométrique déjà explicité dans \hspace{3em} l'exemple précédent.

de $Z_j$, on a $Z'_j\quad Z'_i/N_{ij}$, où $Z'_i$ est l'image inverse de $Z'_j$ sur $Z_i$). Notons que les critères signalés dans l'exemple 1 pour que le foncteur $\underline{r}:\mathrm{Rev}_{\underline{R}}(S)\qquad(\mathrm{Sch})_{/S}$ soit fidèle resp. pleinement fidèle, s'étendent de façon immédiate au cas envisagé ici.

\uline{Exemple} 3 Supposons que $\underline{R}=(\underline{C},\underline{s})$ satisfasse aux conditions suivantes: pour tout $S'\quad\mathrm{Ob}\,\underline{\mathrm{Et}}_S$, et $M\quad\mathrm{Ob}\,\underline{C}_{S'}$, le groupe des automorphismes $G$ de $M$ est fini, et pour toute flèche $M'\quad M$ de $\underline{C}_{S'}$, il existe un unique morphisme de groupe $G'=\mathrm{Aut}(M')$ dans $G=\mathrm{Aut}(M)$ tel que $M'\quad M$ soit compatible avec les opérations de $G'$, de plus si $H$ est son noyau, le morphisme $\underline{s}(M')\quad\underline{s}(M)$ induit un isomorphisme $\underline{s}(M')/H\quad\underline{s}(M)$. \deleted{\unsure{Enfin}} On suppose \deleted{la catégorie $\underline{C}_{S'}$ filtrante (i.e. de pour deux objets $M$, $M'$ de \unsure{ladite}, il existe un objet $M''$ et des flèches $M\quad M''$, $M'\quad M''$),}\inserted{\deleted{que} de plus que pour deux morphismes de $M'$ dans $\unsure{M^{\deleted{\prime}}|M'|M}$, \deleted{les morphismes correspondants de $G'$ dans $G$ sont \unsure{conjugués,|conjugués} \unsure{déduit} l'un de l'autre par composition avec un}} l'un est déduit de l'autre par composition avec un élément de $G$. Enfin, on suppose qu'étant donnés deux objets $M$; $M'$ de $\underline{C}_{S'}$, il existe une famille surjective de morphismes étales $S'_i\quad S'$ telle que pour tout $i$, $M_{S'}$ et $M'_{S_i}$ soient majorés par un objet $M''_i$ de $\underline{C}_{S'_i}$. Utilisant la théorie de la descente, on voit sous ces conditions que \unsure{la|le} foncteur réalisation géométrique \deleted{de} $\mathrm{Rev}_{\underline{R}}(S)\qquad(\mathrm{Sch})_{/S}$ est défini, et à valeurs dans les \deleted{schémas finis sur} revêtements de $S$ (en supposant que les $\underline{s}(M)$ sont des revêtements, bien sûr; pour pouvoir appliquer la descente, il suffirait même que les $\underline{s}(M)$ soient affines sur $S'$). Ici encore, on a des critères pour que ce foncteur soit fidèle resp. pleinement fidèle: il suffit que pour tout $S'\quad\mathrm{Ob}\,\underline{\mathrm{Et}}_S$ et tout $M\quad\mathrm{Ob}\,\underline{C}_{S'}$, la réalisation géométrique $\underline{s}(M)$, munie du groupe d'automorphismes $G=\mathrm{Aut}(M)$, satisfasse aux conditions \inserted{correspondantes} envisagées dans l'exemple 1 pour $(Z,G)$.

\uline{Cas particuliers du type kummérien}: a) Supposons données des sections
% archive sheet number: 19
